% Author: Akshat (condensed from report/section_data_tokenization_retrieval.md, part D)
\section{Data, splits and metric}
\label{sec:data}

\paragraph{Task and decision rule.} Each S2/S3 record is linked to at most one S1 entity; 26\% are linked to
none. Every model uses the same \textbf{decision rule}: a record is assigned to its best-scoring S1 candidate
if the score reaches a threshold $\tau$, and to no S1 otherwise. This is the mention-ranking formulation of
coreference with a dummy antecedent: each record either picks its best antecedent or
starts a new cluster.

\paragraph{Metric.} For every S1 entity we compare the set of S2/S3 records predicted for it with its true
set, compute \fhalf, and average over all S1 entities (macro \fhalf). Singletons, i.e.\ S1 entities with no
true match, are included and score 1 only if nothing is linked to them. Our implementation scores 1.0 on the
ground truth, 0.05585 on an empty prediction (the singleton share), and 0.87114 when the last true record of
every S1 entity is dropped (P 1.0, R 0.66).

\paragraph{Splits.} The official test set has no labels, so all results come from splits of the labelled
training data. The unit is the \textbf{S1 entity}, split 70/10/20 into train/val/test; linked S2/S3 records
follow their S1, so no business is seen in both training and test. Splits are stratified by country $\times$
number of true matches (0 / 1--3 / 4+). Unlinked S2/S3 records (distractors) are spread in the same
proportions by a hash of their id, which keeps the distractor density of the original data. All random
choices use BLAKE2b hashes. We tried CRC32 first, but CRC hashes of one id under different seeds are
correlated.

\paragraph{Regional subset.} Training on the full data (2.2M S1 entities) is beyond free Kaggle T4 GPUs,
so we keep four \emph{whole} states: Oregon and Utah (US), Kerala and Madhya Pradesh (India). Whole states
keep a region's near-duplicate businesses together, so matching stays as hard as on the full data.
States are spelled differently in each source (\emph{TX}/\emph{Texas}). The split script learns the mapping
from address components to states from the ground-truth links: 81 canonical states, 15,095 spellings.
Table~\ref{tab:splits} gives the sizes.

\begin{table}[h]
\centering\small
\begin{tabular}{lrrrrrr}
\toprule
split & S1 (US / India) & singletons & true links & S2 & S3 & unlinked pool \\
\midrule
train & 75,353 (36,238 / 39,115) & 4,169 & 261,183 & 173,057 & 182,320 & 94,194 (26.5\%) \\
val   & 10,764 (5,176 / 5,588)    & 595   & 37,236  & 24,896  & 25,856  & 13,516 (26.6\%) \\
test  & 21,529 (10,353 / 11,176)  & 1,191 & 74,619  & 49,479  & 52,095  & 26,955 (26.5\%) \\
\bottomrule
\end{tabular}
\caption{Splits of the regional subset (\texttt{data/small}).}
\label{tab:splits}
\end{table}

\paragraph{Protocol.} Everything (models, TF-IDF vocabularies, the transliteration lexicon, normalisation
rules) is fitted on train. Hyper-parameters and thresholds are tuned on val. Test is scored once per final
model, and its labels are read only inside the final scoring function.

\paragraph{Limitations of the splits.} (i)~Our splits have about 1.2 unlinked records per S1 entity; on the
official test set we measured about 2.3 during the challenge. Our \fhalf{} is therefore optimistic, mostly in
precision. (ii)~The subset holds about 3.5\% of the S1 entities, so absolute numbers are not those of a
full-data model; the ranking of methods should carry over.
